\section{Software}\label{section:methods:software}

	% ========================================================================
	% SOFTWARE
	%
	% This section documents the analysis software written for this
	% dissertation. Each pipeline gets its own subsection; "Pipeline for SCAPE"
	% is the first.
	%
	% Everything below was reconstructed from the source in
	%   ~/Desktop/Boston_University/Devor_Lab/apical-dendrites-2025
	% specifically code/Preprocessing-STEP1, code/Masks-STEP2, code/Traces-STEP3,
	% the project README.md, and .kiro/steering/pipeline-context.md.
	%
	% Numbers quoted here are the values hardcoded in the scripts at the time of
	% writing. Where a script's configuration varies per dataset, that is stated
	% explicitly. \needsinput{...} blocks mark the places where the code shows
	% WHAT is done but not WHY, and where the physical or biological rationale
	% has to come from you.
	% ========================================================================

	\subsection{Pipeline for SCAPE}\label{section:methods:software:scape}

		Volumetric \acrshort{scape} recordings were analysed with \mavca{}
		(Mask-Assisted Volumetric Calcium Analysis), a modular pipeline written
		for this dissertation in Python~3.11.
		The pipeline takes a 4D ($T \times Z \times Y \times X$) image stack and
		returns background-corrected $\Delta F/F_0$ time series for individual
		dendritic segments, together with the segment masks, curation
		provenance, and behavioural alignment needed to interpret them.

		The design problem \mavca{} addresses is that apical dendrites are thin,
		sparsely active, and densely interleaved.
		A dendritic segment occupies only a few voxels across, fires in a small
		fraction of frames, and lies close enough to its neighbours that a
		generous region of interest will average several segments together.
		Conventional soma-oriented segmentation, which looks for compact
		high-contrast blobs that are active often enough to be found by
		correlation or matrix factorisation, does not recover such structures.
		\mavca{} therefore inverts the usual order: it first finds the moments
		when something was active, uses only those moments to define spatial
		masks, and only then extracts time series for the whole recording.

		\needsinput{%
			The paragraph above states the motivation as I inferred it from the
			code's structure and from the note in
			\texttt{pipeline-context.md} that apical dendrites are only 2--4
			voxels wide. Please confirm or replace it with your own framing, and
			say explicitly which established tools you tried or considered first
			(Suite2p, CaImAn, CNMF-E, Ilastik, manual ImageJ tracing) and how
			they failed on this data. A committee will ask ``why not Suite2p?''
			and the answer is much stronger coming from you.%
		}

		\subsubsection{Data organisation and acquisition parameters}

			Each recording is identified by a triplet of date, mouse, and run,
			and lives in a fixed directory tree
			(\texttt{scape-data/}$\langle$date$\rangle$\texttt{/}$\langle$mouse$\rangle$\texttt{/}$\langle$run$\rangle$) with
			one subdirectory per pipeline stage: \texttt{raw/},
			\texttt{preprocessed/}, \texttt{labelmaps/},
			\texttt{labelmaps\_split/}, \texttt{labelmaps\_curated\_dynamic/},
			\texttt{traces/}, \texttt{trigger/}, \texttt{behavior/} and
			\texttt{overlays/}.
			Stages read from the preceding directory and never overwrite their
			input, so any stage can be re-run without repeating the ones before
			it, and the intermediate state of a recording is visible on disk.

			Voxels are strongly anisotropic, \SI{3.9}{\micro\meter} along $Z$ by
			\SI{1.0}{\micro\meter} along $Y$ by \SI{1.2}{\micro\meter} along $X$,
			so all volume and distance calculations in the pipeline are performed
			in micrometres rather than voxels.
			In the \acrshort{scape} geometry $Y$ corresponds to cortical depth,
			with the pial surface at low $Y$, while $X$ and $Z$ are lateral;
			maximum-intensity projections used for display and for mask review
			are therefore taken as $Y \times X$ views.
			$Z$ is the galvanometer sweep axis, acquired as 30 sequential steps
			at approximately \SI{6.7}{\milli\second} per plane, with depth ($Y$)
			and $X$ captured simultaneously within each $Z$ snapshot.

			This scan order has a direct analytical consequence that the pipeline
			exploits: because a full $Y$ column is acquired in a single snapshot,
			latency measured \emph{along cortical depth} is free of scan-order
			artefact, whereas apparent recruitment \emph{along $Z$} is
			confounded by the sweep and requires subtracting
			$z_\mathrm{centroid} \times \SI{6.7}{\milli\second}$ before
			interpretation.

			Volume rate depends on the acquisition configuration.
			Single-channel 30-plane sessions were acquired at \SI{5}{\hertz};
			deeper dual-channel sessions ($Z = 69$ and $Z = 55$) were acquired at
			\SI{6}{\hertz}.
			Frequency-independent measures (pairwise correlation, effective
			dimensionality, participation fractions, event-class proportions) are
			unaffected by this difference, but every quantity expressed per unit
			time -- event rate, inter-event interval, event duration -- must be
			computed with the volume rate of its own session.

			\needsinput{%
				Two things here need you.
				First, the optical basis of the voxel geometry: why $Z$ is
				sampled roughly four times more coarsely than $Y$ and $X$, what
				sets the axial resolution of the \acrshort{scape} system, and
				how the \acrshort{psf} compares to the
				\SIlist{3.9;1.0;1.2}{\micro\meter} sampling. If the sampling is
				coarser than the \acrshort{psf} along $Z$, say so, because it
				bounds what within-mask propagation could ever be resolved.
				Second, please give the excitation wavelength, power at the
				sample, objective and \acrshort{na}, detection filters, and the
				imaged volume in microns, plus the depth range spanned by $Y$
				relative to the pial surface. I did not find these in the code
				and will not invent them; they belong in
				\cref{section:methods:optics} with a cross-reference here.%
			}

			\needsinput{%
				The raw files carry names of the form
				\texttt{runB\_run1\_rbp4\_132\_phpeb-reslice-bin.tif}, each with
				a companion \texttt{-reslice-bin-dff.tif}, and the project
				\texttt{README} refers to a \texttt{binimagej} variant. This
				implies a reslice, a spatial binning, and a $\Delta F/F$
				computation performed \emph{before} Python, presumably in ImageJ
				or in the \acrshort{scape} acquisition software. Please describe
				that upstream step: the binning factor and the axes it was
				applied to, what the reslice reorients, how the pre-computed
				$\Delta F/F$ stack was calculated, and whether any of it is
				irreversible. It matters because the voxel size quoted above is
				presumably post-binning, and because stage~M1.5 scores frames on
				that externally-generated $\Delta F/F$ stack.%
			}

			\needsinput{%
				\texttt{find\_events\_m1.py} crops three voxels from the deep
				end of $Y$ (\texttt{Y\_CROP = 3}) on every recording, and the
				number of initial seconds discarded is set per session
				(\SI{14}{\second} for 2026-05-08, \SI{12}{\second} for
				2026-05-12 and 2026-04-16). What artefact does the $Y$ crop
				remove, and what is discarded at the start of a recording --
				shutter or laser settling, indicator bleaching, or the animal
				settling after head fixation? The per-session values need a
				stated rule rather than a list.%
			}

		\subsubsection{Stage M1: event detection}

			The first stage identifies the frames in which any activity occurred,
			so that later segmentation can be restricted to them.

			Each voxel's time series is first rescaled to the unit interval by
			its own minimum and maximum over the recording, which equalises the
			contribution of bright and dim voxels.
			The spatial mean of each frame is then subtracted, removing
			field-wide fluctuations shared across the whole volume, and the
			result is smoothed with an isotropic Gaussian kernel of
			$\sigma = 1$~voxel.
			A single scalar activity score per frame is taken as the maximum over
			the smoothed volume, so that a bright event confined to one dendrite
			is not diluted by the inactive remainder of the field.

			Slow drift in that score is removed multiplicatively.
			A rolling baseline $B$ is computed as the 10th percentile within a
			300-frame window, and the score $F$ is converted to
			$F_\mathrm{corr} = F/B - 1$.
			Robust $z$-scores are then formed with \emph{local} rather than
			global statistics: within a 100-frame window the median and the
			median absolute deviation are computed, and
			\begin{equation}\label{equation:methods:m1z}
				z(t) = \frac{F_\mathrm{corr}(t) - \mathrm{median}_\mathrm{loc}(t)}
				            {1.4826 \cdot \mathrm{MAD}_\mathrm{loc}(t)},
			\end{equation}
			where the factor 1.4826 rescales the median absolute deviation to a
			Gaussian-equivalent standard deviation.
			Using local statistics keeps the detection threshold meaningful when
			the noise level itself changes over the course of a recording.

			Events are extracted from $z(t)$ by hysteresis rather than by a
			single threshold: an event opens when $z$ rises above $0.3$ and
			closes only when $z$ falls below $-0.5$, and is retained only if it
			lasts at least three frames.
			Retained events separated by no more than two frames are merged, and
			each is exported as a mini-stack extending five frames either side of
			its extent, together with a table of event onsets, offsets and
			durations in both frames and seconds.
			The asymmetric threshold pair means a segment that is unambiguously
			recruited is followed until it has clearly returned to baseline,
			rather than being fragmented into several short detections whenever
			$z$ dips momentarily.

			\needsinput{%
				The thresholds ($0.3$ opening, $-0.5$ closing, three-frame
				minimum, two-frame merge gap, five-frame crop radius) are the
				values in the script, but the code does not record how they were
				chosen. Were they set by sweeping against manually scored
				events, by visual inspection, or by a target false-positive
				rate? A sentence on provenance, and ideally a supplementary
				sensitivity check showing that the event count is stable across a
				range of thresholds, would make this stage defensible. Note also
				that the three-frame minimum is \SI{600}{\milli\second} at
				\SI{5}{\hertz} but \SI{500}{\milli\second} at \SI{6}{\hertz} --
				worth stating whether the criterion is meant to be in frames or
				in time.%
			}

			\needsinput{%
				Subtracting the spatial mean of each frame is appropriate for
				\emph{detecting} events, but \texttt{pipeline-context.md}
				records that it makes the resulting stack unusable for
				neuropil-contamination tests, because removing a per-frame
				spatial mean forces in-mask and out-of-mask signals to
				anticorrelate at $r \approx -0.97$ as an arithmetic artefact.
				This is an important and easily-misread point. Please confirm
				my reading, and confirm that every contamination or
				neuropil analysis reported in the results chapters was run on
				raw rather than mean-subtracted data.%
			}

		\subsubsection{Stage M1.5: best-frame selection}

			Segmenting from every frame of an event is wasteful and adds noise,
			so this stage selects the few frames in which dendritic structure is
			clearest.

			Frames are scored for \emph{sparse brightness}.
			Each candidate frame is reduced to a maximum-intensity projection
			over the fifteen most superficial $Z$ planes, high-pass filtered by
			subtracting a heavily blurred copy of itself ($\sigma = 15$~voxels)
			and clamping negative values to zero, and scored as the difference
			between its 99.9th and 60th intensity percentiles.
			A frame containing a few bright thin structures against a dark
			background scores highly; one that is uniformly bright, whether from
			background or from a motion artefact, does not.

			Before scoring, the lowest-scoring frame of the event is taken as a
			static-background estimate and $0.8$ times it is subtracted from
			every frame, with negatives clamped to zero.
			Scaling by less than unity deliberately under-subtracts, so that
			dimmer dendrites survive the operation.
			Up to five frames are then retained per event, constrained to be at
			least three frames apart, and any frame scoring below half the best
			frame's score is discarded.
			Each retained frame is written both as a 3D volume, for use by the
			segmentation stage, and as a contrast-stretched projection for visual
			inspection.

			\needsinput{%
				Why the fifteen most superficial $Z$ planes? Given
				\SI{3.9}{\micro\meter} $Z$ sampling this is a
				\SI{58.5}{\micro\meter} slab out of a 30-plane volume, so
				scoring deliberately ignores half the acquired depth. Is that
				because the apical tuft sits in that slab, because the deeper
				planes are dominated by out-of-focus light, or for another
				reason? The same question applies to stage M2's exclusion of the
				most superficial 18\% of $Y$ from thresholding, which the code
				comments attribute to ``surface''. Naming the physical structure
				being excluded in each case -- and confirming the two are
				independent decisions on different axes -- would close an
				obvious gap.%
			}

		\subsubsection{Stage M2: automatic segmentation}

			Masks are proposed automatically from the best frames and event
			crops. Each crop is enhanced by subtracting a heavily blurred copy
			from a lightly smoothed one and clamping to positive values, which
			suppresses slowly-varying background while preserving thin
			structures. The most superficial 18\% of $Y$ is excluded from
			threshold estimation.

			Detection is a seeded hysteresis in space.
			Within each of the four strongest temporal peaks in the crop
			(constrained to be at least five frames apart), a temporal
			maximum-intensity projection over a ten-frame window is thresholded
			twice: at the 99.7th percentile to give high-confidence seeds, and at
			the 96.5th percentile to give a permissive candidate set.
			Connected components of the candidate set are retained only if they
			contain a seed.
			The result is that a structure must be unambiguously bright
			\emph{somewhere} to be admitted at all, but once admitted it is
			followed out to its dimmer extent -- which is what recovers a thin
			dendrite whose distal portion is faint.

			Three optional enhancements address specific failure modes and were
			enabled in the configuration used for the recordings analysed here.
			A Sato vesselness filter over scales
			$\sigma \in \{0.5, 1, 2, 3, 4, 5\}$, thresholded at the 97th
			percentile for seeds and the 95th for candidates, responds to
			elongated tubular structures and improves recovery of trunks.
			An intensity-growth step dilates accepted masks three times, keeping
			only voxels above the 96th intensity percentile, which fills trunk
			bodies that the high-pass enhancement had hollowed out.
			A minimum-intensity projection thresholded at the 2nd percentile
			detects dark outlines, with a \SI{5000}{\micro\meter\cubed} floor.
			Seeds from stage M1.5 are unioned with the automatically detected
			candidates.

			Proposed masks are then cleaned and filtered: per-slice
			morphological opening and closing (kernel sizes 1 and 5), removal of
			slices containing fewer than seven pixels, 3D connected-component
			labelling, a volume floor of \SI{6000}{\micro\meter\cubed} with no
			upper limit, relaxed geometry filters requiring a minimum $Y$ extent
			and elongation, and finally deduplication of any pair whose Dice
			overlap exceeds $0.70$. At most 255 masks are kept per recording.

			\needsinput{%
				Three points need your physical judgement.
				(i)~The \SI{6000}{\micro\meter\cubed} volume floor -- what
				dendrite length does that correspond to at these voxel
				dimensions, and what is the smallest structure you intend to
				accept? Expressing the floor as an equivalent length would be
				far more interpretable to a reader than a volume.
				(ii)~The Sato scales are in voxel units, which are anisotropic,
				so a given $\sigma$ means \SI{3.9}{\micro\meter} along $Z$ but
				\SI{1.0}{\micro\meter} along $Y$. Please confirm this was
				intended, and if the filter was applied to a projection or to
				resampled isotropic data, say which.
				(iii)~The minimum-intensity projection step detects \emph{dark}
				structures, which the code comments call ``dark shadow masks
				from moving trunks''. What is physically dark here -- shadowing
				by a vessel, a trunk displaced by motion, or something else?
				This is unusual enough that a reader will stop at it, and it
				needs one clear sentence.%
			}

		\subsubsection{Stage M2b: splitting merged masks}

			Automatic segmentation sometimes fuses neighbouring dendrites into a
			single mask. Two remedies are available: an automatic
			distance-transform watershed that cuts at geometric constrictions,
			and an interactive \napari{} tool providing polyline cut and erase
			operations.

			One decision here is deliberate and worth stating plainly.
			Roughly \SI{48}{\percent} of split masks consist of a genuine
			dendritic body joined to a stray blob by a neck one voxel wide, and
			such necks are separated by a single erosion.
			Automating that erosion was rejected, because apical dendrites are
			themselves only two to four voxels across and the same operation that
			removes a spurious neck can amputate a real dendrite.
			Neck separation is therefore performed under visual inspection rather
			than offline, accepting slower curation in exchange for not silently
			truncating the structures the study is about.

		\subsubsection{Stage M3: mask curation}

			Every mask is inspected individually in \napari{} against its
			neighbours and either kept, deleted, or edited.
			Keyboard-driven operations support merging, subtraction, retaining
			only the largest connected component, erasing the component under
			the cursor, and lasso selection applied across all $Z$ planes.

			Curation is logged rather than merely applied.
			Deletions are appended to \texttt{curation\_log.csv} and merges to
			\texttt{merge\_log.csv}, masks are renumbered sequentially after any
			edit, and destructive operations write a timestamped backup of the
			mask directory first.
			A classifier trained on accumulated keep/delete decisions
			(\texttt{mask\_classifier.py}) is available to pre-screen proposals.

			\needsinput{%
				This subsection needs the quantities that make manual curation
				auditable: how many masks were proposed and how many survived,
				per recording and in total; the main reasons for rejection; and
				whether curation was blind to condition. If any subset was
				curated twice, by you on separate occasions or by a second
				person, report the agreement -- it is the single most effective
				answer to the objection that hand-curated masks are subjective.
				Please also say whether the classifier was actually used for the
				results reported in this dissertation or was only prototyped; the
				code exists but I cannot tell from it which datasets, if any, it
				was applied to, and the distinction matters for reproducibility.%
			}

		\subsubsection{Stage M4: trace extraction}

			Time series are extracted from curated masks for the full recording.

			A per-voxel baseline volume $F_0$ is computed as the 10th percentile
			over time, after discarding the initial seconds of the recording, and
			relative fluorescence is formed voxel-wise as
			$\Delta F/F_0 = (F - F_0)/F_0$.
			A low percentile rather than a mean is used so that the baseline
			reflects the quiescent state even for segments that are active for a
			substantial fraction of the recording.

			Each mask is then decomposed into a core and a surrounding shell.
			The core is the mask eroded by a unit ball, which discards the
			boundary voxels most likely to contain background; for segments too
			thin to survive erosion the mask is used unchanged.
			The shell is the difference between dilations of the mask by radius~3
			and radius~2, giving a ring one voxel thick and separated from the
			mask by a one-voxel gap.
			The reported trace is the difference of their spatial means,
			\begin{equation}\label{equation:methods:m4}
				\frac{\Delta F}{F_0}(t) =
				\left\langle \frac{\Delta F}{F_0} \right\rangle_{\mathrm{core}}(t)
				-
				\left\langle \frac{\Delta F}{F_0} \right\rangle_{\mathrm{shell}}(t).
			\end{equation}
			The gap exists so that light from the dendrite itself, spread into
			neighbouring voxels by the point spread function, does not enter the
			background estimate and cause the signal to be subtracted from
			itself.

			Residual motion artefacts appear as sharp negative excursions and are
			repaired rather than removed: frames with $\Delta F/F_0 < -0.5$ or a
			frame-to-frame decrement below $-0.3$ are flagged, the flag is
			widened by one frame on each side, and flagged samples are replaced
			by linear interpolation from their neighbours.
			Traces are finally smoothed along time with a Gaussian of
			$\sigma = 0.5$~frames.
			Volumes are streamed in chunks of 120 frames so that memory use is
			independent of recording length.

			This local background subtraction is not a cosmetic step, and its
			omission produced a substantial artefact during development.
			In one field of view, traces extracted from mask cores alone gave a
			mean pairwise correlation of $0.91$ across dendrites, implying
			near-total co-activation; the identical masks with shell subtraction
			applied gave $0.06$.
			Independent testing showed that a field-wide signal is present
			throughout the imaged volume in every recording, correlating at
			$r \approx 0.99$ between in-mask and out-of-mask regions of raw data,
			so the discriminating step is whether that shared signal is removed
			locally. Restricting the core more aggressively did not substitute
			for shell subtraction.

			\needsinput{%
				The field-wide signal removed by shell subtraction needs a
				physical identification from you -- haemodynamic absorption,
				out-of-focus fluorescence from neuropil, laser power
				fluctuation, brain motion, or a combination. It is present in
				all recordings at $r \approx 0.99$, so it is a real property of
				the preparation rather than a per-session defect, and the
				results chapters lean on the claim that subtracting it is
				legitimate rather than that it discards signal. Please also
				comment on whether the one-voxel gap is sufficient given the
				axial \acrshort{psf}: along $Y$ one voxel is
				\SI{1}{\micro\meter}, but a radius-3 ball spans
				\SI{11.7}{\micro\meter} along $Z$, so the shell is far from
				isotropic in physical space.%
			}

			\needsinput{%
				Two smaller choices to justify. First, the artefact rule
				replaces flagged frames by interpolation rather than marking
				them missing, which means an interpolated sample is
				indistinguishable from a measured one downstream; please state
				what fraction of frames this affected and confirm that event
				detection on these traces cannot fire on interpolated segments.
				Second, smoothing with $\sigma = 0.5$~frames is mild but is
				applied before any event statistics are computed -- worth
				stating whether the reported event amplitudes and durations come
				from smoothed or unsmoothed traces.%
			}

		\subsubsection{Concatenation across runs and shared masks}

			Consecutive runs frequently image the same field of view, and
			analyses that follow a given dendrite across runs require that mask
			identities refer to the same physical structure.
			Masks are curated once and applied to sibling runs by reference,
			rather than curated independently per run.
			Which runs legitimately share masks was verified by hashing the
			curated mask sets and confirming byte-identity, because a matching
			field of view is not sufficient -- a merge applied to one run and not
			its sibling silently breaks the correspondence.
			Cross-run analyses keyed on mask identity are restricted to verified
			identical-mask groups.

			For concatenated extraction, $F_0$ is recomputed independently for
			each run rather than pooled across the concatenation, so that a shift
			in baseline between runs is not read as a signal change.
			Slow bleaching is removed by fitting an exponential.
			Run boundaries are retained and marked in all concatenated outputs.

			\needsinput{%
				The bleaching correction needs its model stated: a single
				exponential fitted to what -- each trace individually, the field
				mean, or the raw baseline -- and over what interval. Please also
				say whether bleaching was corrected in the primary
				single-run analyses or only in concatenated ones, since the
				scripts differ in this respect, and whether an exponential was
				adequate or merely convenient.%
			}

		\subsubsection{Grouping masks into branches}

			Because segmentation operates on activity rather than on morphology,
			one dendritic branch may be represented by several masks.
			Candidate groupings are proposed by combining spatial and temporal
			evidence: pairs of masks whose Jaccard overlap exceeds $0.05$ and
			whose traces correlate above $0.3$ during active periods, with at
			least ten active frames available, are linked, and connected sets of
			linked masks are treated as one branch.
			Spatial adjacency graphs built with a \SI{4}{\micro\meter} radius
			were also used to reconstruct branch families and assign a branch
			order and a distance to the deepest mask.

			This reconstruction is fragile.
			Its output is sensitive to the adjacency threshold, and quantities
			derived from it were found to be unstable: the relationship between
			branch order and event rate changed sign between fields of view and
			did not survive a change of adjacency criterion.
			Accordingly, morphological claims in this dissertation are restricted
			to measures that do not depend on family reconstruction -- cortical
			depth, mask volume, mask length -- and event classes defined by
			reconstructed subtrees are reported as exploratory only.

			\needsinput{%
				This is a limitation stated on the strength of the code's own
				notes, and it should be yours to phrase rather than mine, since
				it constrains what the results chapters may claim. If you have
				since improved the reconstruction, or have anatomical ground
				truth for even a handful of branches, that would change the
				conclusion and should replace this paragraph.%
			}

		\subsubsection{Alignment with behavioural signals}

			Behavioural measurements are acquired on separate clocks and must be
			brought onto the imaging timebase before any coupling is computed.
			Accelerometer samples carry a timestamp already aligned to the
			\acrshort{scape} clock and are used directly, taking the raw
			magnitude channel.
			Pupil and whisker traces are acquired by a separate camera and
			require an explicit offset, obtained by comparing the rising edges of
			the camera exposure trigger with those of the imaging trigger
			recorded in the same file; this offset is applied before any
			cropping. The imaging trigger is a sparse start trigger rather than a
			per-frame pulse, so frame times are reconstructed from the start edge
			and the volume rate.

			Coupling between calcium and behaviour is quantified with
			event-triggered averages rather than whole-trace correlation.
			The signals differ in characteristic frequency, and a calcium
			transient and a behavioural variable can be reliably locked in time
			while correlating with opposite sign over the full trace -- in one
			recording a sharp calcium rise preceded an eye movement accompanied
			by a \emph{decrease} in pupil diameter. A whole-trace Pearson
			coefficient would misrepresent such a relationship.

			\needsinput{%
				Please supply the behavioural acquisition details the code does
				not contain: camera model and frame rate (the analysis assumes
				\SI{10}{\hertz} for pupil), how pupil diameter was extracted and
				with what software, where the accelerometer was mounted and what
				its magnitude channel physically measures, and how whisking was
				quantified. Also state the measured magnitude and jitter of the
				camera-to-imaging offset, since the correction is only as good
				as that estimate, and note that two dual-channel sessions have
				pupil and whisker recordings but no accelerometer file.%
			}

		\subsubsection{Implementation and availability}

			\mavca{} is implemented in Python~3.11 and depends on
			NumPy~2.2.5, SciPy~1.15.2, scikit-image~0.25.2,
			scikit-learn~1.6.1, tifffile~2025.3.30, zarr~3.0.7,
			pandas~2.2.3, Matplotlib~3.10.1 and \napari{}~0.6.0 for interactive
			curation; the pinned environment is recorded in
			\texttt{requirements.txt} and checked by \texttt{verify\_env.py}.
			Stage parameters are declared as named constants at the head of each
			script, and the values used for this dissertation are those listed in
			\cref{table:methods:mavca}.

			\needsinput{%
				Decide where the code will live and add it here plus in
				\cref{section:methods:availability}: a public repository with a
				release tag and DOI (Zenodo will mint one from a GitHub
				release), and a licence. Two practical caveats to resolve before
				release, both visible in the code -- absolute paths under
				\texttt{/Users/daria/\ldots} are hardcoded in every script, and
				dataset selection is by editing constants rather than by
				argument, so a reader could not run it as-is. Neither affects
				the validity of the analysis, but both affect the
				reproducibility claim, and it is better to state the limitation
				than to have a reader discover it.%
			}

			% ----------------------------------------------------------------
			% Parameter table. Values transcribed from the scripts as of the
			% time of writing; update if you retune anything.
			% ----------------------------------------------------------------
			\begin{table}[htbp]
				\centering
				\caption[\mavca{} pipeline stages and parameters]{
					Stages of the \mavca{} pipeline for \acrshort{scape} data,
					with the principal parameter values used in this
					dissertation. Percentiles are of intensity unless stated
					otherwise; volumes are in physical units using a voxel size
					of \SIlist{3.9;1.0;1.2}{\micro\meter} along $Z$, $Y$ and
					$X$.
				}\label{table:methods:mavca}
				\begin{singlespace}
				\begin{tabular}{@{}llp{0.44\textwidth}@{}}
					\toprule
					Stage & Purpose & Principal parameters \\
					\midrule
					M1   & Event detection
					     & Gaussian $\sigma = 1$ voxel; rolling baseline 10th
					       percentile over 300 frames; local median and
					       \acrshort{mad} over 100 frames; hysteresis
					       $z > 0.3$ to $z < -0.5$; minimum 3 frames; merge gap
					       2 frames; crop radius 5 frames \\
					\addlinespace
					M1.5 & Best-frame selection
					     & Score $=$ $P_{99.9} - P_{60}$ after high-pass
					       $\sigma = 15$; temporal background $0.8\times$
					       quietest frame; top 15 $Z$ planes; up to 5 frames per
					       event, $\geq 3$ frames apart; score floor $0.5\times$
					       best \\
					\addlinespace
					M2   & Automatic segmentation
					     & 4 temporal peaks $\geq 5$ frames apart; 10-frame
					       temporal projection; seed $P_{99.7}$, candidate
					       $P_{96.5}$; Sato scales 0.5--5; growth at $P_{96}$,
					       3 dilations; volume $\geq$
					       \SI{6000}{\micro\meter\cubed}; Dice deduplication
					       $0.70$ \\
					\addlinespace
					M2b  & Splitting merged masks
					     & Distance-transform watershed; interactive polyline
					       cut and erase; neck separation not automated \\
					\addlinespace
					M3   & Curation
					     & Manual, logged to \texttt{curation\_log.csv} and
					       \texttt{merge\_log.csv}; sequential renumbering;
					       timestamped backups \\
					\addlinespace
					M4   & Trace extraction
					     & $F_0 =$ 10th percentile over time; core $=$ erosion
					       by ball(1); shell $=$ dilation ball(3) minus
					       dilation ball(2); artefact threshold
					       $\Delta F/F_0 < -0.5$ or decrement $< -0.3$;
					       smoothing $\sigma = 0.5$ frames \\
					\bottomrule
				\end{tabular}
				\end{singlespace}
			\end{table}
